\section{Pendahuluan}

Kanker muncul ketika sel yang membawa perubahan genetik atau epigenetik
memperoleh kemampuan untuk membelah, bertahan hidup, dan berekspansi secara
tidak normal. Pada tingkat sel, proses ini berkaitan dengan akumulasi mutasi,
kegagalan DNA repair, checkpoint yang tidak cukup menahan pembelahan, dan
apoptosis yang tidak berhasil mengeliminasi sel berisiko
\cite{hanahan2022,alberts2015}. Karena itu, risiko kanker tidak hanya dapat
dibaca sebagai masalah ukuran organisme, tetapi juga sebagai hasil dinamika
populasi sel. Dengan fokus pada pembelahan sel, checkpoint, akumulasi mutasi,
dan eliminasi sel rusak, proyek ini ditempatkan dalam topik siklus sel.

Secara intuitif, organisme yang lebih besar dan hidup lebih lama diperkirakan
memiliki risiko kanker lebih tinggi. Mereka memiliki lebih banyak sel dan lebih
banyak kesempatan pembelahan sepanjang hidup. Namun, Peto's Paradox
menyatakan bahwa risiko kanker antarspesies tidak meningkat secara
proporsional terhadap ukuran tubuh dan umur \cite{peto1977,caulin2011}.
Paradoks ini mengarah pada pertanyaan mekanistik: mekanisme apa yang dapat
mengurangi peluang sel kanker muncul atau lolos dari kontrol organisme besar?

Studi ini menempatkan simulasi sebagai kontribusi utama. Kami membangun model
populasi sel yang merepresentasikan sel sehat, termutasi, prakanker, dan
kanker. Model ini menguji bagaimana pembelahan, mutasi, repair, apoptosis,
jumlah mutasi yang dibutuhkan untuk transformasi, serta kontrol imun
memengaruhi cancer escape. Dalam laporan ini, cancer escape tidak berarti
kanker klinis pada organisme nyata, melainkan kondisi operasional dalam simulasi
ketika jumlah unit sel kanker melampaui ambang kontrol imun.

Analisis dataset mamalia dari Vincze dkk. \cite{vincze2022} digunakan sebagai
konteks empiris. Dataset tersebut memuat adult cancer mortality risk (CMR),
incidence of cancer mortality (ICM), massa tubuh, adult life expectancy, jumlah
individu, jumlah catatan postmortem, dan jumlah kasus neoplasia. Analisis ini
menunjukkan apakah pola data nyata mendukung motivasi dasar simulasi: ukuran
tubuh dan umur saja tidak cukup untuk menjelaskan variasi risiko kanker.

Pertanyaan penelitian yang diajukan adalah:
\begin{enumerate}
    \item Bagaimana model populasi sel sederhana dapat menerjemahkan mutasi
          multi-hit, repair, apoptosis, dan kontrol imun menjadi cancer escape?
    \item Apakah skenario organisme besar dan berumur panjang tanpa mekanisme
          kompensasi menghasilkan escape probability lebih tinggi daripada
          skenario dengan perlindungan biologis?
    \item Bagaimana hasil simulasi tersebut dibaca bersama analisis empiris CMR
          lintas mamalia dari dataset Vincze dkk.?
\end{enumerate}

Kontribusi studi ini adalah kerangka eksplorasi yang menghubungkan teori
biologi kanker dengan pola Peto's Paradox. Model tidak dimaksudkan untuk
memprediksi risiko kanker spesies tertentu. Tujuannya adalah menunjukkan secara
terkendali bagaimana mekanisme yang termotivasi secara biologis dapat
menurunkan cancer escape pada kondisi ukuran dan umur yang lebih besar.
